\documentclass[a4paper,11pt]{article}
\usepackage{exam-style-842-unit}
\examinehead{842 信息化概论 · 第06单元练习册}{练习一 参考答案与解析}
\begin{document}

\unitanswerscover{842 信息化概论 · 单元练习册}{第06单元参考答案}{练习一（物联网概念与体系架构）的答案与解析；每道判断题均给出改正后的正确表述；四层架构与业界三层的对应关系见相关解析。}

\answerpart{一、单项选择题（10 题，每题 3 分，共 30 分）}

\q B。\begin{analysis}这是教材 7.1.1 给出的物联网定义：物联网即通过射频识别（RFID）、红外感应器、全球定位系统、激光扫描器等信息传感设备，按约定的协议，把任何物品与互联网相连接，进行信息交换和通信，以实现智能化识别、定位、跟踪、监控和管理的一种网络（教材 7.1.1，第 113 页）。A 传感网是由大量智能传感器组成的信息感知网络，只是物联网的组成部分；C 泛在网是物联网的发展远景；D 互联网是物联网的核心和基础，但定义强调“物物相连”。\end{analysis}

\q A。\begin{analysis}教材 7.1.1 指出，物联网的实质包含两层含义：其一，物联网的核心和基础仍然是互联网，是在互联网基础上延伸和扩展的网络；其二，用户端延伸和扩展到了任何物品与物品之间，进行信息交换和通信（教材 7.1.1，第 113 页）。B 传感网是物联网的前身和组成部分；C 泛在网是发展远景；D 云计算属于支撑层的技术手段，不是物联网的核心基础。\end{analysis}

\q B。\begin{analysis}教材图 7.1 和 7.1.2 把物联网的基本构成自下而上分为感知层、网络层、支撑层、应用层四层（教材 7.1.2，第 114—115 页）。A 是业界通行的三层划分，教材在 7.1.2 之外另设支撑层；C、D 层次顺序与教材不符。本题的正确答案即任务中强调的“教材四层”，与业界三层不同，务必按教材作答。\end{analysis}

\q C。\begin{analysis}支撑层主要是在高性能网络计算环境下，将网络内大量或海量信息资源通过计算整合成一个可互联互通的大型智能网络，为上层的服务管理和大规模行业应用建立一个高效、可靠和可信的网络计算超级平台，其主要设备包括大型计算机群、海量网络存储设备、云计算设备等（教材 7.1.2，第 115 页）。A 感知层负责采集；B 网络层负责接入传输；D 应用层是最高层，面向行业建立应用平台。\end{analysis}

\q B。\begin{analysis}网络层主要通过现有互联网（IPv4/IPv6 网络）、移动通信网（GSM、WCDMA、CDMA、无线接入网、无线局域网等）、卫星通信网等基础网络设施，对来自感知层的信息进行接入和传输，其主要设备是能接入各种异构网络的网关等（教材 7.1.2，第 115 页）。A 感知层主要设备是 RFID、传感器、GPS、GIS 等；C 支撑层主要设备是大型计算机群、海量存储、云计算设备；D 应用层主要是用户界面显示设备及管理设备。\end{analysis}

\q B。\begin{analysis}射频识别技术是物联网中让物品“开口说话”的关键技术：RFID 标签上存储着规范而具有互用性的信息，可通过无线数据通信网把它们自动采集到中央信息系统，实现物品的识别（教材 7.1.2，第 115 页）。A 二维码也是感知技术，但教材明确把“让物品开口说话”的表述给 RFID；C GPS 用于移动物体定位；D 多媒体采集处理用于视频、音频、图像信息的采集。\end{analysis}

\q A。\begin{analysis}感知层主要用于采集物理世界中发生的物理事件和数据，即利用 RFID、传感器、GPS、GIS 等随时随地获取物体的信息，实现物理世界的全面感知（教材 7.1.2，第 114 页）。B 负责传输；C 负责计算与存储；D 面向行业应用。\end{analysis}

\q A。\begin{analysis}教材 7.1.1 对 3 个术语作了对比：传感网是由大量智能传感器组成的信息感知网络，是物联网的前身，现在是物联网的组成部分（教材 7.1.1，第 114 页）。B、C 把二者关系说反或等同；D 泛在网才是物联网的发展远景，传感网不是。\end{analysis}

\q A。\begin{analysis}2009 年 1 月，IBM 首席执行官彭明盛提出“智慧地球”构想，其中物联网为“智慧地球”不可或缺的重要部分，即互联网加物联网等于智慧地球（教材 7.1.5，第 119 页）。英特尔是 2002 年建立世界上第一个无线葡萄园的公司（教材 7.1.3，第 116 页）；微软与比尔·盖茨是 1995 年《未来之路》提及物联网概念；谷歌与本题无关。\end{analysis}

\q C。\begin{analysis}1999 年，EPC global 的 Auto-ID 中心认为，物联网是成千上万的物品采用无线方式接入 Internet 网络，这是物联网概念的真正提出（教材 7.1.5，第 118 页）。A 1995 年比尔·盖茨在《未来之路》中只是提及物联网概念；B 1998 年 MIT 的 Kevin Ashton 把 RFID 技术与传感器应用于日常物品形成“物联网”，早于概念的真正提出；D 2005 年 ITU 报告把物联网概念拓展到了传感网络。\end{analysis}

\answerpart{二、多项选择题（4 题，每题 4 分，共 16 分）}

\q ABC。\begin{analysis}教材 7.1.1 列出物联网的 3 个基本特征：互联网特征（需要联网的物品能够互联互通形成网络）、识别与通信特征（具备自动识别与物物通信能力，利用 RFID、传感器、二维码等获取物体信息，实现全面感知）、智能处理（利用云计算、模糊识别等智能计算技术对海量数据进行分析处理并实施智能控制）。D 错误：物联网强调低功耗、绿色节能，“高能耗传输”不是其特征。（教材 7.1.1，第 113—114 页）\end{analysis}

\q ABC。\begin{analysis}教材 7.1.1 明确指出：物联网是互联网的延伸和扩展，互联网是物联网的核心和基础；用户端延伸和扩展到了任何物品与物品之间，进行信息交换和通信。A、B、C 均符合教材表述。D 错误：物联网不是独立于互联网的新网络，而是建立在互联网基础之上的延伸网络。（教材 7.1.1，第 113 页）\end{analysis}

\q ABC。\begin{analysis}感知层利用 RFID、传感器、GPS、GIS 等随时随地获取物体的信息，实现物理世界的全面感知（教材 7.1.2，第 114 页）。A、B、C 均属于感知层的信息获取技术。D 错误：异构网络网关是网络层的主要设备，负责接入和传输，不属于感知技术。\end{analysis}

\q ABC。\begin{analysis}教材 7.1.5 指出，欧洲智能系统集成技术平台组织（EPoSS）在《Internet of Things in 2020》中预测，物联网 2020 年之后实现全智能化；美国 FORRESTER 预测到 2020 年物物互联业务与人际通信业务之比将达 30:1，物联网被称为下一个万亿美元级的通信业务。A、B、C 正确。D 错误：教材同时指出物联网大规模应用仍面临成本、安全、侵犯隐私 3 个方面的巨大挑战，并非“已不存在挑战”。（教材 7.1.5，第 120—121 页）\end{analysis}

\answerpart{三、判断题（6 题，每题 2 分，共 12 分）}

\q √。\begin{analysis}教材 7.1.1 明确：物联网是物物相连的互联网，其实质的第一层含义就是“物联网的核心和基础仍然是互联网，是在互联网基础上延伸和扩展的网络”。二者不是替代关系。（教材 7.1.1，第 113 页）\end{analysis}

\q ×。改正：教材第 7 章把物联网体系结构划分为感知层、网络层、支撑层、应用层四个层次，其中支撑层是云计算、海量存储的超级平台；业界通行的三层架构把计算与存储支撑并入网络层、应用层来理解，按三层作答也应对应给分。\begin{analysis}教材图 7.1 与 7.1.2 共列出四个层次（教材 7.1.2，第 114—115 页）。这是本单元最需要辨析的一点：三层是业界通行说法，四层是教材口径，答题应以教材四层为主并说明对应关系。\end{analysis}

\q √。\begin{analysis}教材 7.1.2 指出，支撑层的主要设备包括大型计算机群、海量网络存储设备、云计算设备等，其作用是为上层提供网络计算超级平台（教材 7.1.2，第 115 页）。该项表述与教材一致。\end{analysis}

\q ×。改正：感知层负责采集物理世界中发生的物理事件和数据；网络层负责对来自感知层的信息进行接入和传输，其主要设备是能接入各种异构网络的网关。\begin{analysis}教材 7.1.2 对两层的分工有明确表述：感知层“主要用于采集物理世界中发生的物理事件和数据”，网络层“对来自感知层的信息进行接入和传输”（教材 7.1.2，第 114—115 页）。把采集职责安到网络层，混淆了四层架构的层间分工。\end{analysis}

\q √。\begin{analysis}教材 7.1.1 的 3 个术语对比指出：传感网是物联网的前身，现在是物联网的组成部分；泛在网是物联网的发展远景，物联网是泛在网的前期（教材 7.1.1，第 114 页）。\end{analysis}

\q ×。改正：物联网的用户端延伸和扩展到了任何物品与物品之间，可以进行物与物、人与物、人与人的信息交换，并不只限于人与人之间的通信。\begin{analysis}教材 7.1.1 的第二层含义明确指出，用户端延伸和扩展到物品与物品之间，可进行物与物、人与物、人与人的信息交换（教材 7.1.1，第 113 页）。只承认人与人通信就退回到互联网的范畴，抹掉了物联网“物物相连”的本质。\end{analysis}

\answerpart{四、配对题（6 对，每对 1 分，共 6 分）}

\q （1）—D\quad （2）—C\quad （3）—E\quad （4）—B\quad （5）—A\quad （6）—F。
\begin{analysis}感知层采集物理事件和数据；网络层以能接入各种异构网络的网关为主要设备；支撑层由大型计算机群、海量存储和云计算设备构成，是网络计算超级平台；应用层面向绿色农业、智能交通等行业建立应用平台；RFID 是让物品“开口说话”的关键技术；GPS 是全球定位系统，用于采集移动物体信息。（对应关系见教材 7.1.2，第 114—115 页）\end{analysis}

\answerpart{五、简答题（3 题，每题 8 分，共 24 分）}

\q 参考答案要点：
（1）定义（3 分）：物联网即通过射频识别（RFID）、红外感应器、全球定位系统、激光扫描器等信息传感设备，按约定的协议，把任何物品与互联网相连接，进行信息交换和通信，以实现智能化识别、定位、跟踪、监控和管理的一种网络。（教材 7.1.1，第 113 页）
（2）两层含义（2 分）：一是物联网的核心和基础仍然是互联网，是在互联网基础上延伸和扩展的网络；二是用户端延伸和扩展到任何物品与物品之间，进行信息交换和通信，可进行物与物、人与物、人与人的信息交换。
（3）基本特征（3 分）：互联网特征，需要联网的物品一定要能够互联互通形成网络；识别与通信特征，物品要具备自动识别与物物通信能力，利用 RFID、传感器、二维码等随时随地获取物体信息，实现全面感知；智能处理，利用云计算、模糊识别等智能计算技术对海量数据进行分析处理，对物体实施智能化控制。
（4）补充（可不答）：从技术层面看，物联网是物体经智能感应装置、传输网络到达信息处理中心，实现自动化信息交互与处理的智能网络；从应用层面看，是把所有物体连入网络并与互联网结合，实现人类社会与物理系统的整合。
（5）常见失分：只写定义不写两层含义；把“互联网特征”漏掉；把传感网、泛在网的定义混入物联网特征。

\q 参考答案要点：
（1）感知层（2 分）：采集物理世界中发生的物理事件和数据，利用 RFID、传感器、GPS、GIS 等随时随地获取物体信息，实现物理世界的全面感知；主要设备/技术包括 RFID 标签与读写器、各类传感器、GPS、GIS、二维码等。
（2）网络层（2 分）：通过现有互联网（IPv4/IPv6）、移动通信网（GSM、WCDMA、CDMA、无线接入网、无线局域网等）、卫星通信网等基础网络设施，对来自感知层的信息进行接入和传输；主要设备是能接入各种异构网络的网关。
（3）支撑层（2 分）：在高性能网络计算环境下，把网络内大量或海量信息资源通过计算整合成一个可互联互通的大型智能网络，为上层的服务管理和大规模行业应用建立高效、可靠、可信的网络计算超级平台；主要设备包括大型计算机群、海量网络存储设备、云计算设备等。
（4）应用层（2 分）：物联网系统结构的最高层，包括各类用户界面显示设备及其他管理设备；面向各类行业建立实际应用的管理平台和运行平台，并集成相关内容服务，如绿色农业、公共安全、远程医疗、智能交通等，一般以综合管理中心的形式出现。
（5）答题提示：若按业界通行的感知层、网络层、应用层三层作答，应说明教材把云计算、海量存储等计算支撑单列为支撑层，位于网络层与应用层之间；历年真题按三层作答亦应给分。（教材 7.1.2，第 114—115 页；真题 2022 年简答“农业物联网体系结构”即考此点）

\q 参考答案要点：
（1）物联网与互联网（3 分）：互联网是全球范围的信息交互网络，是物联网的核心和基础；物联网是互联网的延伸和扩展，把用户端从计算机等终端扩展到物品与物品之间，实现物物、人物、人人之间的信息交换。二者不是替代关系，物联网没有脱离互联网另建一张网。
（2）物联网与传感网（2 分）：传感网是由大量智能传感器组成的信息感知网络，侧重信息的感知与采集，是物联网的前身，现在是物联网的组成部分；物联网在传感网的基础上增加了网络传输、计算支撑和行业应用，范围更大。
（3）物联网与泛在网（2 分）：泛在网即无所不在的网络社会，强调任何人、任何物、任何时间、任何地点的连接，是物联网的发展远景；物联网是泛在网的前期，是实现泛在网的基础阶段。
（4）四者关系链（1 分）：传感网（感知）$\rightarrow$ 物联网（物物互联）$\rightarrow$ 泛在网（无所不在），互联网始终是物联网的核心和基础网络。
（5）常见失分：把物联网说成“独立于互联网的新网络”；把传感网与物联网等同；把泛在网与物联网的先后关系写反。（教材 7.1.1，第 113—114 页）

\answerpart{六、综合题（1 题，共 12 分）}

\q 参考答案要点：
（1）四层归位（6 分，每层 1.5 分）：
感知层——土壤温湿度传感器、小型气象站、视频监控、北斗终端、稻米包装上的二维码，负责采集土壤、气象、影像、位置和标识信息；
网络层——4G 移动通信网与光纤宽带，以及接入异构网络的网关，负责把感知层数据接入并传输到云平台；
支撑层——农业云平台及其大型计算机群、海量网络存储设备和云计算设备，负责海量数据的存储、计算与整合，向上提供网络计算超级平台；
应用层——手机应用、管理平台与溯源查询界面，面向农户提供灌溉建议、面向消费者提供溯源查询。
支撑层不能省略（2 分）：材料中的“农业云平台”承担的正是海量数据存储与计算，教材把这类计算支撑单列为支撑层；若只写三层，就说不清数据存在哪里、灌溉建议由谁计算，也丢失了教材四层架构的完整性与“超级平台”的定位。
（2）三层对应（2 分）：业界三层为感知层、网络层、应用层。对应关系是：感知层相同；教材网络层与支撑层合并理解为业界的网络层（或把计算存储支撑视为网络层的延伸能力）；应用层相同。答题时应先按教材四层作答，再补一句“支撑层相当于业界三层中网络层之上的计算与存储支撑，部分资料将其并入网络层”，并注明历年真题按三层作答亦给分。
（3）基本特征（2 分）：
互联网特征——传感器、气象站、北斗终端、云平台和手机应用都接入网络，互联互通；
识别与通信特征——二维码标识稻米身份，传感器与北斗终端自动采集并回传数据，实现全面感知；
智能处理——云平台对海量数据进行分析，向农户推送灌溉建议、向消费者提供溯源信息，属于智能计算与智能控制。
（4）结论（1 分）：该园区是“感知—传输—支撑—应用”四层协同的农业物联网系统，四层缺一不可；答题时点明教材四层与业界三层的对应关系，可覆盖 2022 年“农业物联网体系结构”、2024 年“物联网定义、构成和关键技术”等真题的采分点。（教材 7.1.2、7.1.3，第 114—117 页）

\end{document}
